\documentclass[10pt,a4paper]{amsart}
\usepackage{amsmath,amssymb,amsthm,graphicx}
% Retained minimal preamble for the Stacks Project chapter varieties.tex.
% Verbatim excerpt of the chapter preamble: Xy-pic setup (native preamble lines 14-18), the native theorem
% declarations with the mathematical macros (native preamble lines 154-217), and the project own no-op
% definitions of the reference/slogan/history display environments (native wrapper.tex lines 1-3) that the
% counted statement uses for its slogan. Copied mechanically; nothing added, removed or redefined.
% For commutative diagrams we use Xy-pic
\usepackage[all]{xy}

% We use 2cell for 2-commutative diagrams.
\xyoption{2cell}
% Theorem environments.
%
\theoremstyle{plain}
\newtheorem{theorem}[subsection]{Theorem}
\newtheorem{proposition}[subsection]{Proposition}
\newtheorem{lemma}[subsection]{Lemma}

\theoremstyle{definition}
\newtheorem{definition}[subsection]{Definition}
\newtheorem{example}[subsection]{Example}
\newtheorem{exercise}[subsection]{Exercise}
\newtheorem{situation}[subsection]{Situation}

\theoremstyle{remark}
\newtheorem{remark}[subsection]{Remark}
\newtheorem{remarks}[subsection]{Remarks}

\numberwithin{equation}{subsection}

% Macros
%
\def\lim{\mathop{\mathrm{lim}}\nolimits}
\def\colim{\mathop{\mathrm{colim}}\nolimits}
\def\Spec{\mathop{\mathrm{Spec}}}
\def\Hom{\mathop{\mathrm{Hom}}\nolimits}
\def\Ext{\mathop{\mathrm{Ext}}\nolimits}
\def\SheafHom{\mathop{\mathcal{H}\!\mathit{om}}\nolimits}
\def\SheafExt{\mathop{\mathcal{E}\!\mathit{xt}}\nolimits}
\def\Sch{\mathit{Sch}}
\def\Mor{\mathop{\mathrm{Mor}}\nolimits}
\def\Ob{\mathop{\mathrm{Ob}}\nolimits}
\def\Sh{\mathop{\mathit{Sh}}\nolimits}
\def\NL{\mathop{N\!L}\nolimits}
\def\CH{\mathop{\mathrm{CH}}\nolimits}
\def\proetale{{pro\text{-}\acute{e}tale}}
\def\etale{{\acute{e}tale}}
\def\QCoh{\mathit{QCoh}}
\def\Ker{\mathop{\mathrm{Ker}}}
\def\Im{\mathop{\mathrm{Im}}}
\def\Coker{\mathop{\mathrm{Coker}}}
\def\Coim{\mathop{\mathrm{Coim}}}

% Boxtimes
%
\DeclareMathSymbol{\boxtimes}{\mathbin}{AMSa}{"02}

%
% Macros for moduli stacks/spaces
%
\def\QCohstack{\mathcal{QC}\!\mathit{oh}}
\def\Cohstack{\mathcal{C}\!\mathit{oh}}
\def\Spacesstack{\mathcal{S}\!\mathit{paces}}
\def\Quotfunctor{\mathrm{Quot}}
\def\Hilbfunctor{\mathrm{Hilb}}
\def\Curvesstack{\mathcal{C}\!\mathit{urves}}
\def\Polarizedstack{\mathcal{P}\!\mathit{olarized}}
\def\Complexesstack{\mathcal{C}\!\mathit{omplexes}}
% \Pic is the operator that assigns to X its picard group, usage \Pic(X)
% \Picardstack_{X/B} denotes the Picard stack of X over B
% \Picardfunctor_{X/B} denotes the Picard functor of X over B
\def\Pic{\mathop{\mathrm{Pic}}\nolimits}
\def\Picardstack{\mathcal{P}\!\mathit{ic}}
\def\Picardfunctor{\mathrm{Pic}}
\def\Deformationcategory{\mathcal{D}\!\mathit{ef}}
\newenvironment{reference}{}{}
\newenvironment{slogan}{}{}
\newenvironment{history}{}{}
\title{Separation of a k-rational point from finitely many points by a section of a very ample invertible sheaf}
\author{The Stacks Project Authors}
\date{stacks-project:varieties.tex:local-snapshot-20261011; GNU FDL 1.2 or later}
\begin{document}
\maketitle
\noindent\textit{Editorial note.} All delivered mathematical text is the author's own, copied from the pinned Stacks Project chapter file varieties.tex; no statement or proof step was rewritten, reordered or removed.\par\medskip
\noindent\textit{Editorial note.} The author disposes of the case n = 0 with the single word 'trivial' and proceeds; that sentence is reproduced as written and no case is dropped.\par\medskip
\noindent\textit{Editorial note.} The preamble is the chapter's retained minimal preamble (sources/stacks-varieties/retained-preamble.tex lines 1-77), a verbatim excerpt of the native preamble: Xy-pic setup, the native theorem declarations with the mathematical macros, and the project's no-op display environments for slogan/reference/history.\par\medskip
\noindent\textit{Editorial note.} The label lemma-very-ample-vanish-at-point is defined only in varieties.tex, so no cross-chapter restatement had to be weighed for this card.\par\medskip
\noindent\textit{Allowed standard prerequisites (editorial).}\begin{itemize}
\item Very ample invertible sheaf on a finite type k-scheme. Same chapter, the definition of a very ample invertible sheaf as one whose sections give a closed immersion into projective space. Supplies the reduction of the statement to the case X = P\textasciicircum{}r\_k with L = O\_X(1).
\item Homogeneous coordinate ring and Proj of P\textasciicircum{}r\_k. Same chapter, P\textasciicircum{}r\_k = Proj(k[T\_0, ..., T\_r]) and the correspondence between points and homogeneous prime ideals. Supplies the identification of the sections vanishing at x with a codimension one subspace W of V = kT\_0 + ... + kT\_r.
\item A finite dimensional vector space over an infinite field is not a finite union of proper subspaces. Standard linear algebra, used without a pointer at proof lines 4255-4256. Supplies the conclusion that W is not the union of the finitely many W cap p\_i.
\end{itemize}\medskip
\begin{lemma}
\label{lemma-very-ample-vanish-at-point}
Let $k$ be an infinite field. Let $X$ be a scheme of finite type over $k$.
Let $\mathcal{L}$ be a very ample invertible sheaf on $X$.
Let $n \geq 0$ and $x, x_1, \ldots, x_n \in X$ be points with
$x$ a $k$-rational point, i.e., $\kappa(x) = k$, and
$x \not = x_i$ for $i = 1, \ldots, n$.
Then there exists an $s \in H^0(X, \mathcal{L})$ which vanishes at
$x$ but not at $x_i$.
\end{lemma}

\begin{proof}
If $n = 0$ the result is trivial, hence we assume $n > 0$.
By definition of a very ample invertible sheaf, the lemma immediately
reduces to the case where $X = \mathbf{P}^r_k$ for some $r > 0$
and $\mathcal{L} = \mathcal{O}_X(1)$. Write
$\mathbf{P}^r_k = \text{Proj}(k[T_0, \ldots, T_r])$.
Set $V = H^0(X, \mathcal{L}) = kT_0 \oplus \ldots \oplus kT_r$.
Since $x$ is a $k$-rational point, we see that the set
$s \in V$ which vanish at $x$ is a codimension $1$ subspace
$W \subset V$ and that $W$ generates the homogeneous prime
ideal corresponding to $x$. Since $x_i \not = x$ the corresponding
homogeneous prime $\mathfrak p_i \subset k[T_0, \ldots, T_r]$ does
not contain $W$. Since $k$ is infinite, we then see that
$W \not = \bigcup W \cap \mathfrak q_i$ and the proof is complete.
\end{proof}

\end{document}
